\subsection{Research objectives}
\label{sec:objectives}

The key objectives which framed the direction of the research activities centered around delivering a novel FLISR solution, leveraging modern technologies
which can operate both centrally and in a distributed fashion at the network edge, and to test, evaluate and demonstrate the operation and effectiveness of the solution through
the use of a simulation platform, utilising a number of grid topologies to assess the FLISR algorithms effectiveness against a variety of network configurations located throughout Ireland.
\\
The main output of the research is gauging how the developed FLISR solution may reduce CML in real-world operation, to that end, the following questions guided the research activities.

\begin{itemize}
  \item \textbf{R1 - How effective will the outputs of the FLISR solution be in terms of reducing DSO fault KPIs?}
\end{itemize}

\noindent
This question is crucial in regards to evaluating the value of the research activities and FLISR solution implemented, both compared to alternative approaches to the FLISR concept
undertaken in other research, and in relation to the tested grid topologies and extrapolated to the Irish energy grid as a whole.

\begin{itemize}
  \item \textbf{R2 - Given a similar set of conditions, how will the results produced by the centralised FLISR algorithm compare to the distributed FLISR algorithms operating locally at the edge of the network?}
\end{itemize}

\noindent
A key feature of the distributed FLISR solution is to be in itself fault-tolerant, and therefore able to operate and produce results locally in the event communications with the central,
``master'' FLISR algorithm is unavailable. Given the limited view of the network granted to the FLISR algorithms operating locally at the edge,
in comparison to the overall system awareness of the central FLISR algorithm, will the results produced by both be similar and, crucially, valid.

\begin{itemize}
  \item \textbf{R3 - Given a diverse set of network topologies, will the FLISR solution output valid results for a variety of grid configurations?}
\end{itemize}

\noindent
For the FLISR solution to be effective, a key requirement will be to operate and produce valid results given a wide set of grid topologies with varying network configurations in terms
of structure, number of nodes and device types.

\clearpage

% subsection objectives (end)
